\documentclass[12pt]{article}
\usepackage[utf8]{inputenc}
\usepackage[T1]{fontenc}
\usepackage{hyperref}
\title{Prai Language Resources for K-12 Curriculum Development}
\author{Compiled by Claude Sonnet 5.5 (claude-sonnet-5-5) for Andr\'as Kornai}
\date{}
\begin{document}
\maketitle

\section{Introduction}
Prai (Phai, ISO 639-3 prt) is an Austroasiatic language of Thailand and Laos, closely related to Mal (mlf); both are sometimes called Thin or Htin, which makes sources easy to confuse. No school materials were found, and nothing found indicates that it is taught in schools.

\section{Primary Educational Resources}
None found.

\section{Linguistic and Reference Materials}
\begin{itemize}
\item Wikipedia, Prai language (phonological inventory): \url{https://en.wikipedia.org/wiki/Prai_language}
\item Wikipedia, Mal language (close relative; do not conflate): \url{https://en.wikipedia.org/wiki/Mal_language}
\item CRDO/Cocoon language page for prt, with a bibliography of older work (Filbeck 1978, Jordan manuscripts, a 2004 sociolinguistic survey, a Mahidol MA thesis on Pray grammar, and a Nan lexical geography dissertation): \url{https://cocoon.huma-num.fr/exist/crdo/lexvo.xql?code=prt}
\item 2025 sociolinguistic survey of the ``Lua Bo Kluea'' community in Nan Province: \url{https://so07.tci-thaijo.org/index.php/DJIIH/article/download/7217/5998/84584}
\end{itemize}

\section{Religious and Community Resources}
\begin{itemize}
\item Wikipedia, Prai people: \url{https://en.wikipedia.org/wiki/Prai_people}
\item Joshua Project profiles for the Prai in Thailand and Laos: \url{https://joshuaproject.net/index.php/people_groups/14377/TH} and \url{https://joshuaproject.net/index.php/people_groups/14377/LA}
\end{itemize}

\section{Digital and Audio Resources}
None found.

\section{Curriculum Development Context}
Documentation is sparse and mostly sociolinguistic or descriptive. No orthography, primers or school programme were found, so any curriculum would need new materials developed with Prai speakers. The Thai and Lao national curricula are the likely school context.
\end{document}